\section{Polyhedral blocks coupled through a fixed core}\label{app:fixed-core}

This appendix develops a related elimination principle for a single-level
problem that is linear in many small polyhedral blocks and polynomial in
a fixed-dimensional core. It both supplies a broader structural comparison
and makes witness recovery explicit. Its convex combinations preserve
polyhedral feasibility and linear measurements; they are not combinations
of globally optimal responses of a nonconvex follower.

Fix $r,k\ge0$ and $d\ge1$. Let $C\subset\R^r$ be compact and given by an explicit
rational quantifier-free polynomial formula. For $b=1,\ldots,B$, let
\[
 P_b(x)=\{y\in\R^{d_b}:G_b(x)y\le h_b(x)\},\qquad 1\le d_b\le d,
\]
with polynomial rational data and explicit finite rational coordinate
bounds among its rows. The local sets may be empty or lower-dimensional.
Let $A_b(x)\in\Q[x]^{k\times d_b}$, $b(x)\in\Q[x]^k$,
$c_b(x)\in\Q[x]^{d_b}$ and $c_0(x)\in\Q[x]$. Consider
\begin{equation}\label{eq:fixed-core-problem}
 \begin{aligned}
 \min_{x,z}\quad&c_0(x)+\sum_{b=1}^B c_b(x)\trans z_b,\\
 \text{subject to}\quad&x\in C,\quad z_b\in P_b(x)\quad(b=1,\ldots,B),\\
 &\sum_{b=1}^B A_b(x)z_b=b(x).
 \end{aligned}
\end{equation}
Let $L$ include all explicit data and let $\delta\ge2$ bound actual
input degree. In contrast with a fixed-degree restriction, the following
statement tracks the numerical degree throughout its construction.

\begin{theorem}[Fixed-core polyhedral optimization]\label{thm:fixed-core}
For fixed $r,d,k$, feasibility and global optimization of
\eqref{eq:fixed-core-problem} are computable in time polynomial in
$(L,\delta)$. If feasible, an optimum is attained and an optimizer with
all block coordinates and its value can be returned in a single
real-algebraic field with polynomial total encoding length. Neither the
number of blocks nor their row counts is fixed. A fixed number of
aggregate inequalities is also permitted if finite bounds for their
slacks are supplied.
\end{theorem}

\subsection{Local vertices and support directions}\label{app:local-support}

For block $b$ enumerate all $d_b$-row subsets $I$. Put
\[
 \Delta_{bI}(x)=\det (G_b(x))_I,\qquad
 p_{bI}(x)=\operatorname{adj}((G_b(x))_I)(h_b(x))_I.
\]
Where $\Delta_{bI}\ne0$, the candidate is
$p_{bI}/\Delta_{bI}$. Equivalently, use the positive-denominator form
\begin{equation}\label{eq:core-vertex-formula}
 D_{bI}=\Delta_{bI}^2,\qquad Y_{bI}=\Delta_{bI}p_{bI},\qquad
                          y_{bI}=Y_{bI}/D_{bI}.
\end{equation}
Its validity predicate is
\begin{equation}\label{eq:core-vertex-validity}
 D_{bI}>0,\qquad G_b(x)Y_{bI}(x)\le h_b(x)D_{bI}(x).
\end{equation}
All these tests are polynomial of degree $O(d\delta)$.

Every vertex of $P_b(x)$ has $d_b$ independent active rows. Otherwise
there is a nonzero direction orthogonal to all active normals; a small
step in either direction keeps the active rows equal and preserves all
finitely many strict inactive rows, contradicting extremality. A
nonempty bounded polyhedron has a vertex, and every linear functional
has a maximizing vertex. For completeness, a maximizing face is nonempty
and compact; successively minimize the coordinates within it to produce
a singleton face, which is a vertex of the original polyhedron.
These arguments also apply to lower-dimensional sets and singletons.
Thus \eqref{eq:core-vertex-formula} lists enough candidates to detect
nonemptiness and maximize every linear functional.

Append the objective as a measurement. With $h=k+1$, define
\begin{equation}\label{eq:core-images}
 W_b(x)=\begin{pmatrix}A_b(x)\\c_b(x)\trans\end{pmatrix},\quad
 T_b(x)=W_b(x)P_b(x),\quad
 t(x,\eta)=\begin{pmatrix}b(x)\\\eta-c_0(x)\end{pmatrix}.
\end{equation}
For a direction $\lambda\in\R^h$, the support score of a valid
vertex is
\[
 \lambda\trans W_b(x)y_{bI}(x)=n_{bI}(x,\lambda)/D_{bI}(x),\qquad
 n_{bI}=\lambda\trans W_bY_{bI}.
\]
Collect the denominator and validity polynomials, and, for every pair
of candidates in one block, the comparison polynomial
\begin{equation}\label{eq:core-score-comparison}
                     n_{bI}D_{bJ}-n_{bJ}D_{bI}.
\end{equation}
The denominators of valid candidates are positive, so these signs
correctly order their scores. There are polynomially many polynomials
for fixed $d$, all in the $r+h$ coordinates $(x,\lambda)$, of degree
$O(d\delta)$ and polynomial coefficient bit length.

Enumerate their realizable sign conditions, remembering identically zero
polynomials and including zero signs. Each sign condition determines the
valid vertices and a support-maximizing vertex for every nonempty block;
choose the first maximizer in a fixed ordering. Discard a condition if
some block has no valid candidate. The retained list $\Sigma$ is
polynomial in size and can be constructed in polynomial bit time for
fixed $(r,d,k)$. Disconnected sign realizations cause no problem, since
both validity and score order are determined by signs.

\subsection{Exact membership from the support inequalities}\label{app:support-membership}

For $\sigma\in\Sigma$, let $I_b(\sigma)$ be its selected local
index and abbreviate the corresponding denominator and numerator by
$D_{b\sigma},n_{b\sigma}$. Put
\begin{equation}\label{eq:core-cleared-support}
 \begin{split}
 H_\sigma(x)&=\prod_bD_{b\sigma}(x),\\
 J_\sigma(x,\eta,\lambda)&=
 \lambda\trans t(x,\eta)H_\sigma(x)
             -\sum_b n_{b\sigma}(x,\lambda)\prod_{a\ne b}D_{a\sigma}(x).
 \end{split}
\end{equation}
On that sign condition $H_\sigma>0$ and $J_\sigma\le0$ states the
support inequality
\[
 \lambda\trans t(x,\eta)\le
                     \sum_b\max_{y\in P_b(x)}\lambda\trans W_b(x)y.
\]
Writing $S_\sigma(x,\lambda)$ for the sign conjunction, define
\begin{equation}\label{eq:core-support-predicate}
 \Phi(x,\eta,\lambda)=\bigvee_{\sigma\in\Sigma}
                  [S_\sigma(x,\lambda)\wedge J_\sigma(x,\eta,\lambda)\le0].
\end{equation}

\begin{lemma}[Support membership]\label{lem:core-support-membership}
For $x\in C$, there exist blocks satisfying
\eqref{eq:fixed-core-problem} with objective exactly $\eta$ if and only if
\begin{equation}\label{eq:core-membership}
                           \forall\lambda\in\R^h\ \Phi(x,\eta,\lambda).
\end{equation}
\end{lemma}
\begin{proof}
If a block is empty, no sign condition realized at $(x,\lambda)$ is
retained, for any $\lambda$. Thus \eqref{eq:core-membership} is false.
Otherwise the Minkowski sum
$T(x)=\sum_bT_b(x)=\{\sum_b t_b:t_b\in T_b(x)\}$ is compact and
convex. Its support function is the sum of the support functions of its
summands: a maximizing choice can be made independently in every block.
The formula \eqref{eq:core-support-predicate} therefore gives its exact
support inequality in every direction.

A point $t$ in a compact convex set $T$ satisfies all its support
inequalities. Conversely, if $t\notin T$, choose a nearest point
$p\in T$ and put $\lambda=t-p\ne0$. The derivative of squared
distance along the segment from $p$ to any $u\in T$ gives
$(t-p)\trans(u-p)\le0$. Hence
$\max_{u\in T}\lambda\trans u\le\lambda\trans p
<\lambda\trans t$, violating a support inequality.
Thus \eqref{eq:core-membership} is equivalent to $t(x,\eta)\in T(x)$.
By \eqref{eq:core-images}, this membership states precisely existence
of blocks satisfying all linking equations and attaining the specified
objective value.
\end{proof}

The degrees in \eqref{eq:core-cleared-support} are
$O((B+1)d\delta)$, which need not be constant. Their ambient dimension
$r+h+1$ is fixed, however, so their expanded monomial count and bit size
are polynomial in $(L,\delta)$ by the product argument in
\cref{lem:candidate-size}. Define
\[
 E(x,\eta)=(x\in C)\wedge\forall\lambda\,\Phi(x,\eta,\lambda).
\]
This fixed-variable formula describes exact attainable core--value
pairs. The full feasible set of \eqref{eq:fixed-core-problem} is closed
and lies in the product of compact $C$ and the finite block boxes.
Therefore the upper minimum is attained whenever feasible. Apply
quantifier elimination and joint sampling to
\begin{equation}\label{eq:core-joint-optimum}
 E(x,\eta)\wedge
       \neg\exists x',\eta'\,[E(x',\eta')\wedge\eta'<\eta].
\end{equation}
This gives an optimal $(x^*,\eta^*)$ in one field
$K=\Q(\alpha)$ of polynomial degree and encoding length.
It remains to recover all block variables. Giving each of them a new
quantified coordinate would destroy the fixed-dimensional argument.

\subsection{Constructive recovery from a few support tuples}\label{app:core-recovery}

The same support enumeration already supplies a polynomial-size set of
full block tuples sufficient for recovery. For each $\sigma\in\Sigma$
test its selected local denominators at $x^*$ and discard the tuple
if any is zero. Evaluate the remaining formulas and retain the tuple
if every selected local vertex is feasible at $x^*$. Denote the retained tuples by $y^1,\ldots,y^M$, and set
\[
              a^j=\sum_b W_b(x^*)y_b^j\in K^h\qquad(j=1,\ldots,M).
\]
There is no need to require that the original sign condition $\sigma$
be realizable at $x^*$ with some direction: feasibility of the tuple
suffices for this filtering step. It never adds a point outside $T(x^*)$.

\begin{lemma}[A small common convex combination recovers every block]
\label{lem:core-recovery}
The retained aggregate points satisfy
$\conv\{a^1,\ldots,a^M\}=T(x^*)$. For every target $t\in K^h\cap T(x^*)$
one can find at most $h+1$ retained tuples and nonnegative weights in $K$
that sum to one and whose weighted aggregate is $t$, in polynomial bit
time in the original input and target encoding. Using the same weights
in every block gives a feasible original tuple with aggregate $t$.
\end{lemma}
\begin{proof}
Every retained tuple is blockwise feasible, so its aggregate belongs to
$T(x^*)$ and the indicated convex hull is contained in $T(x^*)$.
For every real direction $\lambda$, the actual sign condition at
$(x^*,\lambda)$ is retained in $\Sigma$: no block is empty at an
optimal core. Its selected vertices give a feasible tuple that attains
the support of $T(x^*)$ in direction $\lambda$. This tuple survives
the filter. Hence the convex hull has the same support function as
$T(x^*)$. The separation argument in \cref{lem:core-support-membership}
proves equality.

The usual Carath\'eodory reduction gives a convex representation of
a target in a finite convex hull with affinely independent support. To see this directly, start with a
representation with all supported weights positive. If its lifted
columns $(a^j,1)$ are dependent, choose a nonzero dependence $u$.
Its components sum to zero, so it has a positive component. Replace
the weights $\theta_j$ by $\theta_j-su_j$, with
$s=\min_{u_j>0}\theta_j/u_j$. This preserves the target and the sum
of weights, keeps weights nonnegative, and removes at least one support
point. Repeat until the lifted columns are independent. Their number
is at most $h+1$.

For construction, enumerate all subsets of at most $h+1$ retained
points. On each affinely independent subset solve
\[
          \sum_j\theta_j a^j=t,\qquad \sum_j\theta_j=1.
\]
Choose independent rows to solve this fixed-size system, then check
all remaining equations and the nonnegativity of its unique weights.
If a suitable subset exists over the reals, its independent linear
system has its unique solution in $K$, since every coefficient and the
target belong to $K$. Thus the enumeration finds it. There are at most
$\sum_{q=1}^{h+1}\binom Mq$ subsets, polynomial for fixed $h$.
All arithmetic, zero tests, and comparisons take place in the already
sampled ordered field $K$, whose degree and coefficient encodings are
polynomial. Each system has fixed order, so determinant expressions
and rational evaluation give polynomial coefficient lengths for the
weights; summing the resulting fixed number of terms preserves the
same bound for every recovered block coordinate.

Finally, put $z_b=\sum_j\theta_j y_b^j$ for every block. Convexity
of $P_b(x^*)$ makes $z_b$ feasible, and linearity of $W_b(x^*)$ gives
$\sum_bW_b(x^*)z_b=t$. The same weights must be used in all blocks:
this is how the prescribed total measurements are preserved.
\end{proof}

Applying \cref{lem:core-recovery} to
$t=t(x^*,\eta^*)$ recovers the original optimizer in the field $K$,
completing the main claim of \cref{thm:fixed-core}. With no blocks,
the image sum is $\{0\}$ and the empty tuple is its sole support tuple;
all the formulas use empty sums and products with their usual meanings.
Blocks of dimension zero can be removed after testing their scalar rows.
A fixed number of aggregate inequalities can be made equalities using
bounded scalar slack blocks. Supplied slack bounds ensure that the
result stays in the same compact model and do not change the fixed
structural dimensions.

This argument gives polynomial dependence on the numerical input degree
as well as input length: local determinants and pair comparisons have
degree $O(d\delta)$; the global products have degree
$O((B+1)d\delta)$; sign enumeration, elimination and sampling are all
in fixed dimension; and reconstruction adds only fixed-size algebraic
linear systems. Under polynomially bounded degree, this is polynomial
time in $L$.

The support characterization, finite-dimensional convex-hull reduction,
and real algebraic algorithms are established tools. An alternative
recovery would solve a linear program over the sampled common field;
Adler and Beling~\cite{AdlerBeling1994} analyze that route with explicit
dependence on the common extension degree. The constructive recovery
above uses the support tuples already generated and needs no
large-dimensional algebraic LP solve. It depends on the affine block
measurements, including the objective coordinate. It would not justify
convexifying a nonconvex follower reaction set or mixing its tied
responses to create additional follower choices.
